    \documentclass{article}
    \usepackage{amsmath}
	\usepackage{amssymb}
    \usepackage[minted]{tcolorbox}
    \usepackage{xcolor} % Required for bgcolor in minted
    \definecolor{lightgray}{rgb}{0.9,0.9,0.9} % Define lightgray if not already defined
    \usepackage{fontspec} %fontspec is required for code font
    \setmonofont{JuliaMono}[Extension=.ttf, UprightFont=*-Regular, BoldFont=*-Bold, ItalicFont=*-RegularItalic, BoldItalicFont=*-BoldItalic, Contextuals=Alternate, Scale = 0.8]
    \usepackage{tabularx} % For tables that fit page width
    \usepackage{array} % For better column formatting
    \usepackage{booktabs} % For professional looking tables
	\usepackage{fullpage,graphicx,psfrag,amsfonts,verbatim, url}
    \usepackage{parskip}
    
    \title{\textbf{\textsf{\texttt{PEPit.jl} tutorial: Accelerated proximal point for monotone inclusions}}}
	
	\author{John Doe}
    
    \begin{document}
    \maketitle
    
    

Before we do anything let us load the necessary packages!

\begin{tcblisting}{listing only, minted language=julia, minted options={mathescape=true, breaklines}, colback=lightgray, colframe=black, boxrule=0pt, toprule=2pt, bottomrule=2pt, arc=0pt, outer arc=0pt, left=5pt, right=5pt, top=5pt, bottom=5pt}
using PEPit, OrderedCollections, Mosek, MosekTools, OffsetArrays
\end{tcblisting}


Let us start with defining an empty PEP, which we will construct step by step.

\begin{tcblisting}{listing only, minted language=julia, minted options={mathescape=true, breaklines}, colback=lightgray, colframe=black, boxrule=0pt, toprule=2pt, bottomrule=2pt, arc=0pt, outer arc=0pt, left=5pt, right=5pt, top=5pt, bottom=5pt}
problem = PEP()
\end{tcblisting}


Consider the monotone inclusion problem

\begin{equation*}
\mathrm{Find}\, x:\, 0\in Ax,
\end{equation*}

where $A$ is maximally monotone. We denote $J_A = (I + A)^{-1}$ the resolvent of $A$. Let us define this operator class in \texttt{PEPit} next.

\begin{tcblisting}{listing only, minted language=julia, minted options={mathescape=true, breaklines}, colback=lightgray, colframe=black, boxrule=0pt, toprule=2pt, bottomrule=2pt, arc=0pt, outer arc=0pt, left=5pt, right=5pt, top=5pt, bottom=5pt}
A = declare_function!(problem, MonotoneOperator, OrderedDict())
\end{tcblisting}


\textbf{Algorithm}: The algorithm that we want to investigate is the accelerated proximal point is described as follows, for $t \in \{ 0, \dots, n-1\}$

\begin{equation*}
\begin{aligned}
    x_{t+1} &= J_{\alpha A}(y_t), \\
    y_{t+1} &= x_{t+1} + \frac{t}{t+2}(x_{t+1} - x_{t}) - \frac{t}{t+2}(x_t - y_{t-1}),
\end{aligned}
\end{equation*}

where $x_0=y_0=y_{-1}$,  where $\alpha$ is a positive scalar. 

Our goal is to compute the smallest possible $\tau(n, \alpha)$ such that the guarantee

\begin{equation*}
\|x_n - y_n\|^2 \leqslant \tau(n, \alpha) \|x_0 - x_\star\|^2,
\end{equation*}

is valid, where $A x_\star \ni 0$ and $x_0$ is the initial point that we pick.

For our tutorial, let us take $\alpha = 2.0$ and $n=10$.

\begin{tcblisting}{listing only, minted language=julia, minted options={mathescape=true, breaklines}, colback=lightgray, colframe=black, boxrule=0pt, toprule=2pt, bottomrule=2pt, arc=0pt, outer arc=0pt, left=5pt, right=5pt, top=5pt, bottom=5pt}
alpha = 2

n = 10
\end{tcblisting}


We next define the optimal point $x_\star \equiv$ \texttt{xs} that  satisfies $0 \in Ax_\star$. and the initial point $x_0 \equiv $ \texttt{x0},

\begin{tcblisting}{listing only, minted language=julia, minted options={mathescape=true, breaklines}, colback=lightgray, colframe=black, boxrule=0pt, toprule=2pt, bottomrule=2pt, arc=0pt, outer arc=0pt, left=5pt, right=5pt, top=5pt, bottom=5pt}
xs = stationary_point!(A)

x0 = set_initial_point!(problem)
\end{tcblisting}


The next stage is defining the initial condition $\|x_0 - x_\star\|^2$.

\begin{tcblisting}{listing only, minted language=julia, minted options={mathescape=true, breaklines}, colback=lightgray, colframe=black, boxrule=0pt, toprule=2pt, bottomrule=2pt, arc=0pt, outer arc=0pt, left=5pt, right=5pt, top=5pt, bottom=5pt}
set_initial_condition!(problem, (x0 - xs)^2 <= 1)
\end{tcblisting}


Now it is time to implement our algorithm in a loop.

\begin{tcblisting}{listing only, minted language=julia, minted options={mathescape=true, breaklines}, colback=lightgray, colframe=black, boxrule=0pt, toprule=2pt, bottomrule=2pt, arc=0pt, outer arc=0pt, left=5pt, right=5pt, top=5pt, bottom=5pt}
x = OffsetVector(fill(x0, n + 1), 0:n)
y = OffsetVector(fill(x0, n + 1), 0:n)

for i in 0:(n - 2)
    x[i + 1], _, _ = proximal_step!(y[i + 1], A, alpha)
    y[i + 2] = x[i + 1] + i / (i + 2) * (x[i + 1] - x[i]) - i / (i + 2) * (x[i] - y[i])
end
x[n], _, _ = proximal_step!(y[n], A, alpha)
\end{tcblisting}


Not it is time to define our performance metric, which is $\|x_n - y_n\|^2$.

\begin{tcblisting}{listing only, minted language=julia, minted options={mathescape=true, breaklines}, colback=lightgray, colframe=black, boxrule=0pt, toprule=2pt, bottomrule=2pt, arc=0pt, outer arc=0pt, left=5pt, right=5pt, top=5pt, bottom=5pt}
set_performance_metric!(problem, (x[n] - y[n])^2)
\end{tcblisting}


We have everything to build up the PEP. Now time to solve!

\begin{tcblisting}{listing only, minted language=julia, minted options={mathescape=true, breaklines}, colback=lightgray, colframe=black, boxrule=0pt, toprule=2pt, bottomrule=2pt, arc=0pt, outer arc=0pt, left=5pt, right=5pt, top=5pt, bottom=5pt}
τ_PEPit = solve!(problem, solver=Mosek.Optimizer, verbose=true)
\end{tcblisting}


If everything is well, we should get \texttt{τ\_PEPit = 0.01}  for \texttt{n = 10} and \texttt{alpha = 2}.

What does the theory say? A tight theoretical worst-case guarantee can be found in [1, Theorem 4.1],
for $n \geqslant 1$,

\begin{equation*}
\|x_n - y_{n-1}\|^2 \leqslant  \frac{1}{n^2}  \|x_0 - x_\star\|^2.
\end{equation*}

\begin{tcblisting}{listing only, minted language=julia, minted options={mathescape=true, breaklines}, colback=lightgray, colframe=black, boxrule=0pt, toprule=2pt, bottomrule=2pt, arc=0pt, outer arc=0pt, left=5pt, right=5pt, top=5pt, bottom=5pt}
τ_theory = 1 / (n^2)

@info "💻  PEPit guarantee: τ_PEPit = $(round(τ_PEPit, digits=6))"

@info "📚  Theoretical guarantee: τ_theory = $(round(τ_theory, digits=6))"
\end{tcblisting}


We should see: 

\begin{tcblisting}{listing only, minted language=julia, minted options={mathescape=true, breaklines}, colback=lightgray, colframe=black, boxrule=0pt, toprule=2pt, bottomrule=2pt, arc=0pt, outer arc=0pt, left=5pt, right=5pt, top=5pt, bottom=5pt}
PEPit guarantee: τ_PEPit = 0.01

Theoretical guarantee: τ_theory = 0.01
\end{tcblisting}


So the guarantee that we found via \texttt{PEPit} is the same as theoretical guarantee!

\textbf{Reference}:

[1] D. Kim (2021). Accelerated proximal point method for maximally monotone operators.
Mathematical Programming, 1-31.
<https://arxiv.org/pdf/1905.05149v4.pdf>


\end{document}
